\usepackage{amsmath,amsfonts,amssymb}
\usepackage{graphics}
\usepackage{graphicx}
\usepackage{booktabs}
\usepackage{tabularx}
\usepackage{array}
\usepackage{icomma}
\usepackage{tikz, ifthen}
\usetikzlibrary{arrows,shapes,backgrounds,patterns,decorations.pathreplacing,decorations.pathmorphing}

% <Must include this>
% Definitions for coin images
\def\coinheads{\draw (0,0) circle (0.9cm); \draw (0,0) circle (0.72cm); \draw (0,0) node {\Huge\textbf{H}};}
\def\cointails{\draw (0,0) circle (0.9cm); \draw (0,0) circle (0.72cm); \draw (0,0) node {\Huge\textbf{T}};}
% Definitions for dice images
\def\makedieempty#1#2{
  \begin{scope}[#1]
    \pgfsetcornersarced{\pgfpoint{#2}{#2}}
    \draw (-0.5cm, -0.5cm) rectangle (0.5cm, 0.5cm);
  \end{scope}
}

\def\makedieone#1#2{
  \begin{scope}[#1]
    \pgfsetcornersarced{\pgfpoint{#2}{#2}}
    \draw (-0.5cm, -0.5cm) rectangle (0.5cm, 0.5cm);
    \fill[black] (0, 0) circle (0.125);
  \end{scope}
}

\def\makedietwo#1#2{
  \begin{scope}[#1]
    \pgfsetcornersarced{\pgfpoint{#2}{#2}}
    \draw (-0.5cm, -0.5cm) rectangle (0.5cm, 0.5cm);
    \fill[black] (-0.3, -0.3) circle (0.125);
    \fill[black] (0.3, 0.3) circle (0.125);
  \end{scope}
}

\def\makediethree#1#2{
  \begin{scope}[#1]
    \pgfsetcornersarced{\pgfpoint{#2}{#2}}
    \draw (-0.5cm, -0.5cm) rectangle (0.5cm, 0.5cm);
    \fill[black] (-0.3, -0.3) circle (0.125);
    \fill[black] (0, 0) circle (0.125);
    \fill[black] (0.3, 0.3) circle (0.125);
  \end{scope}
}

\def\makediefour#1#2{
  \begin{scope}[#1]
    \pgfsetcornersarced{\pgfpoint{#2}{#2}}
    \draw (-0.5cm, -0.5cm) rectangle (0.5cm, 0.5cm);
    \fill[black] (-0.3, -0.3) circle (0.125);
    \fill[black] (0.3, -0.3) circle (0.125);
    \fill[black] (0.3, 0.3) circle (0.125);
    \fill[black] (-0.3, 0.3) circle (0.125);
  \end{scope}
}

\def\makediefive#1#2{
  \begin{scope}[#1]
    \pgfsetcornersarced{\pgfpoint{#2}{#2}}
    \draw (-0.5cm, -0.5cm) rectangle (0.5cm, 0.5cm);
    \fill[black] (0, 0) circle (0.125);
    \fill[black] (-0.3, -0.3) circle (0.125);
    \fill[black] (0.3, -0.3) circle (0.125);
    \fill[black] (0.3, 0.3) circle (0.125);
    \fill[black] (-0.3, 0.3) circle (0.125);
  \end{scope}
}

\def\makediesix#1#2{
  \begin{scope}[#1]
    \pgfsetcornersarced{\pgfpoint{#2}{#2}}
    \draw (-0.5cm, -0.5cm) rectangle (0.5cm, 0.5cm);
    \fill[black] (-0.3, -0.3) circle (0.125);
    \fill[black] (-0.3,  0.0) circle (0.125);
    \fill[black] (-0.3,  0.3) circle (0.125);
    \fill[black] ( 0.3, -0.3) circle (0.125);
    \fill[black] ( 0.3,  0.0) circle (0.125);
    \fill[black] ( 0.3,  0.3) circle (0.125);
  \end{scope}
}

\def\makediegeneral#1#2#3{
  \ifthenelse{#1=1}{\makedieone{#2}{#3}}{
    \ifthenelse{#1=2}{\makedietwo{#2}{#3}}{
      \ifthenelse{#1=3}{\makediethree{#2}{#3}}{
        \ifthenelse{#1=4}{\makediefour{#2}{#3}}{
          \ifthenelse{#1=5}{\makediefive{#2}{#3}}{
            \ifthenelse{#1=6}{\makediesix{#2}{#3}}{
              \makedieempty{#2}{#3}}}}}}}
}

\def\makedie#1#2{
  \makediegeneral{#1}{#2}{0.175cm}
}

\def\makediesmall#1#2{
  \makediegeneral{#1}{#2}{0.0875cm}
}

% Definitions for Venn images
\def\samplespace{(0,0) rectangle (1,1)}
\def\labelsamplespace{\draw (1,1) node[anchor=south] {$S$}}

\def\circlepartiala{(0.25,0.5) circle (0.2)}
\def\labelpartiala{\draw (0.25,0.7) node[anchor=south] {$A$}}
\def\circlepartialb{(0.6,0.4) circle (0.3)}
\def\labelpartialb{\draw (0.6,0.7) node[anchor=south] {$B$}}

\def\circleseparatea{(0.2,0.3) circle (0.1)}
\def\labelseparatea{\draw (0.23,0.25) node[anchor=north west] {$A$}}
\def\circleseparateb{(0.5,0.6) circle (0.2)}
\def\labelseparateb{\draw (0.5,0.8) node[anchor=south] {$B$}}

\def\circlefulla{(0.5,0.5) circle (0.35)}
\def\labelfulla{\draw (0.75,0.75) node[anchor=south west] {$A$}}
\def\circlefullb{(0.5,0.6) circle (0.15)}
\def\labelfullb{\draw (0.5,0.45) node[anchor=north] {$B$}}

% </Must include this>
   






\newenvironment{toc}%
    {\begin{description}\setlength{\topsep}{0cm}\setlength{\itemsep}{0cm}%%
    \setlength{\parskip}{0cm}\setlength{\parsep}{0cm}%%
    \setlength{\partopsep}{0cm}}
    {\end{description}}

    \newenvironment{indexheading}%
    {\begin{description}\setlength{\topsep}{0cm}\setlength{\itemsep}{0cm}%%
    \setlength{\parskip}{0cm}\setlength{\parsep}{0cm}%%
    \setlength{\partopsep}{0cm}
    \setlength{\labelwidth}{0cm}\setlength{\labelsep}{0cm}%
    \setlength{\leftmargin}{-1cm} \item[] \noindent}
    {\end{description}}


    \usepackage{marginnote}
  
    \usepackage[paper=letter, 
      lmargin=0.1\paperwidth, rmargin=0.25\paperwidth, 
      tmargin=1in, bmargin=1in,
      twoside, centering,
      includehead, marginsep=3mm, marginparwidth=0.225\paperwidth ]{geometry}
   \usepackage{etoolbox} 
    \usepackage{wrapfig}
    \usepackage{ifthen,changepage}      
    \usepackage{layouts}
    \usepackage{multicol}
    \usepackage{amsmath} %for \underset, \overset, and more?
    \usepackage{amssymb} %for set of reals, integers, etc..., \triangleq
    \usepackage{alltt}   %for codeblocks
    \usepackage{url} %for nice url breaks
    \usepackage{graphicx} % for figure support
    \usepackage{epstopdf}
    \usepackage{subfigure}
    \usepackage{tabularx}
    \usepackage{supertabular}
    \usepackage{xtab}
    \usepackage{multirow}
    \usepackage{float}
    \usepackage{ragged2e} 
    \usepackage{array}
    \usepackage{mathrsfs}
    \usepackage{textcomp}
    \usepackage[cjkjis,mathletters,autogenerated,tipa]{ucs}
    \usepackage{amsmath, amsthm, amsfonts, amssymb}
    \usepackage{xcolor}

    \usepackage{pst-all}
    \usepackage{pst-eucl}        %Jo
    \usepackage{pst-poly}        %Jo
    \usepackage{pst-spectra}
    \usepackage{pst-slpe}
    \usepackage{pst-math}

    %% ************* NB ************
    %% The order in which pstricks packages are loaded
    %% matters - so I copied the order from pst-all.sty
    %% and then added the two additional packages at the end.
    %% ************* End NB ************

    %% ************* Packages ************
    \usepackage{pst-circ}
    \usepackage{pst-lens}
    \usepackage{pst-optic}         %Jo
    \usepackage{pstricks-add}         %Jo
    \usepackage{pst-labo}         %Jo  
    \usepackage{pst-3dplot} 
    \usepackage{pst-solides3d}
    \usepackage{auto-pst-pdf}
    \usepackage{mdframed}
    \sffamily
    \pagestyle{headings}

    
    \usepackage[utf8x]{inputenc}
    \usepackage[C40,T1]{fontenc}
    \usepackage[scaled]{helvet}           
    \renewcommand*\familydefault{\sfdefault}
        
    \usepackage{calc}
    \usepackage{ifthen}
    %\usepackage{breqn}
    \usepackage{ulem}
    \normalem

%    \setlength{\emergencystretch}{3em}
    \renewcommand{\labelenumi}{\textbf{\arabic*{enumi}.}}
    \renewcommand{\thesubfigure}{(\alph{subfigure})}
    \renewcommand{\labelitemi}{\ensuremath{\bullet} }
    \renewcommand{\labelitemii}{\ensuremath{\cdot} }
    \usepackage{enumerate}
    \usepackage{enumitem}
    \setcounter{tocdepth}{3}
    \DeclareUnicodeCharacter {9653} {\bigtriangleup}
    \DeclareUnicodeCharacter {183} {\ensuremath{\cdot}}
    \DeclareUnicodeCharacter {785} {\textasciibreve}
    \DeclareUnicodeCharacter {787} {'}
    \DeclareUnicodeCharacter {788} {`}
    \DeclareUnicodeCharacter {789} {'}
    \DeclareUnicodeCharacter {940} {[U+03AC]}
    \DeclareUnicodeCharacter {941} {[U+03AD]}
    \DeclareUnicodeCharacter {942} {[U+03AE]}
    \DeclareUnicodeCharacter {943} {[U+03AF]}
    \DeclareUnicodeCharacter {970} {\"{i}}
    \DeclareUnicodeCharacter {972} {[U+03CC]}
    \DeclareUnicodeCharacter {973} {[U+03CD]}
    \DeclareUnicodeCharacter {974} {[U+03CE]}
    \DeclareUnicodeCharacter {978}  {\ensuremath{\Upsilon}}
    \DeclareUnicodeCharacter {988}  {\ensuremath{\digamma}}
    \DeclareUnicodeCharacter {8260} {/}
    \DeclareUnicodeCharacter {8289} {}
    \DeclareUnicodeCharacter {8290} {}
    \DeclareUnicodeCharacter {8407} {\ensuremath{\rightarrow}}
    \DeclareUnicodeCharacter {8474} {\ensuremath{\mathbb{Q}}}
    \DeclareUnicodeCharacter {8484} {\ensuremath{\mathbb{Z}}}
    \DeclareUnicodeCharacter {8497} {\ensuremath{\mathcal{F}}}
    \DeclareUnicodeCharacter {8519} {\ensuremath{\mathbb {e}}}
    \DeclareUnicodeCharacter {8520} {\ensuremath{\mathbb {i}}}
    \DeclareUnicodeCharacter {8596} {\ensuremath{\leftrightarrow}}
    \DeclareUnicodeCharacter {8614} {\ensuremath{\mapsto}}
    \DeclareUnicodeCharacter {8788} {:=}
    \DeclareUnicodeCharacter {9001} {\ensuremath{\textless}}
    \DeclareUnicodeCharacter {9002} {\ensuremath{\textgreater}}   
    \DeclareUnicodeCharacter {9474} {\ensuremath{\textbar}}
    \DeclareUnicodeCharacter {10003} {\checkmark}
  
    \DeclareUnicodeCharacter {61168}{\ensuremath{\jmath}}
    \DeclareUnicodeCharacter {61237}{\ensuremath{\mathscr{A}}}
    \DeclareUnicodeCharacter {61238}{\ensuremath{\mathscr{C}}}
    \DeclareUnicodeCharacter {61239}{\ensuremath{\mathscr{D}}}
    \DeclareUnicodeCharacter {61240}{\ensuremath{\mathscr{G}}}
    \DeclareUnicodeCharacter {61241}{\ensuremath{\mathscr{J}}}
    \DeclareUnicodeCharacter {61242}{\ensuremath{\mathscr{K}}}
    \DeclareUnicodeCharacter {61243}{\ensuremath{\mathscr{N}}}
    \DeclareUnicodeCharacter {61244}{\ensuremath{\mathscr{O}}}
    \DeclareUnicodeCharacter {61245}{\ensuremath{\mathscr{P}}}
    \DeclareUnicodeCharacter {61246}{\ensuremath{\mathscr{Q}}}
    \DeclareUnicodeCharacter {61247}{\ensuremath{\mathscr{S}}}
    \DeclareUnicodeCharacter {61248}{\ensuremath{\mathscr{T}}}
    \DeclareUnicodeCharacter {61249}{\ensuremath{\mathscr{U}}}
    \DeclareUnicodeCharacter {61250}{\ensuremath{\mathscr{V}}}
    \DeclareUnicodeCharacter {61251}{\ensuremath{\mathscr{W}}}
    \DeclareUnicodeCharacter {61252}{\ensuremath{\mathscr{X}}}
    \DeclareUnicodeCharacter {61253}{\ensuremath{\mathscr{Y}}}
    \DeclareUnicodeCharacter {61254}{\ensuremath{\mathscr{Z}}}
    \DeclareUnicodeCharacter {61327}{\ensuremath{\mathbb {E}}}
    \DeclareUnicodeCharacter {61328}{\ensuremath{\mathbb {F}}}
    \DeclareUnicodeCharacter {62838}{\ensuremath{\longleftarrow}}
    \DeclareUnicodeCharacter {62839}{\ensuremath{\longrightarrow}}
    \DeclareUnicodeCharacter {62843}{\ensuremath{\leftrightarrow}}




\newenvironment{IFact}[1]%
{\vspace{3mm}\rmfamily
 \begin{center}\begin{pspicture}(0,0)(5,0.01)
 \psline[linecolor=red]{-}(0,0)(5,0)
\end{pspicture}\end{center}
%\psshadowbox{
 \begin{quote}%
\rput(-1.5,0){\psfact}\noindent #1%
\end{quote}%
 }%}
{\begin{center}\begin{pspicture}(0,0)(5,0.01)
 \psline[linecolor=red]{-}(0,-0.1)(5,-0.1)
\end{pspicture}\end{center}
\vspace{3mm}\sffamily}%

\newcommand{\Definition}[2]{
\begin{mdframed}[linewidth=2, margin=1cm,linecolor=lightgray,skipbelow=0.5cm,%
skipabove=0.5cm]%
\textsl{DEFINITION: } \hspace{1cm} #1\nopagebreak\par
#2
\end{mdframed}
}

\newcommand{\Identity}[2]{
\vspace{0.5cm}
\begin{mdframed}[linewidth=3pt,linecolor=lightgray,skipbelow=\baselineskip,%
skipabove=\baselineskip]%
\begin{flushleft}
%\uppercase{\textsl{#1}}\nopagebreak
#2
\end{flushleft}
\end{mdframed}
}


\newcommand{\Tip}[1]{
\marginpar{%odd
\begin{minipage}[m]{0.225\paperwidth}
\begin{mdframed}[linecolor=darkgray, linewidth=1.5pt]
Tip
\end{mdframed}
\begin{mdframed}[backgroundcolor=lightgray, linecolor=lightgray,skipabove=-4pt]
\small #1
\end{mdframed}
\end{minipage}
}
}

\newcommand{\Warning}[1]{
\begin{mdframed}[linewidth=4pt,%
linecolor=red%
]
%\begin{mdframed}[linewidth=2pt,linecolor=red]%
\vspace{3mm}
\begin{center}
\noindent%
\textbf{Warning:}\par
#1
\end{center}
\end{mdframed}
\vspace{3mm}}


\newcommand{\Note}[1]{
\marginpar{%odd
\begin{minipage}[h]{0.225\paperwidth}

\begin{mdframed}[linecolor=siyavula-dark-gray, backgroundcolor=siyavula-dark-gray, skipbelow=-pt, margin=10pt]
Note
\end{mdframed}
\begin{mdframed}[backgroundcolor=siyavula-gray, linecolor=siyavula-gray, margin=10pt, skipabove=-4pt]
\small #1
\end{mdframed}
\end{minipage}
}
}

                

    \setlength{\parskip}{2ex}

    \setlength{\parindent}{0pt}

    
	\newenvironment{note}[1]%
	{\begin{list}{}{%
 	\setlength{\labelsep}{0pt}\setlength{\rightmargin}{20pt}%
        {\leftmargin}{20pt}%
	\setlength{\labelwidth}{0pt}\setlength{\listparindent}{0pt}}%
        \item\textbf{\textsc{#1}}}%
	{\end{list}}

  

    \newenvironment{cnxcaption}%
	{\begin{list}{}{%
 	\setlength{\labelsep}{0pt}\setlength{\rightmargin}{25pt}%
        \setlength{\leftmargin}{25pt}%
	\setlength{\labelwidth}{0pt}\setlength{\listparindent}{0pt}}%
        \item}%
	{\end{list}}

        \newenvironment{example}{%
          \noindent
          \begingroup
          \leftskip=20pt\rightskip=\leftskip
        }{%
          \vspace{\rubberspace}\par
          \endgroup
        }

        %\newenvironment{exercise}{\begin{example}}
        %{\end{example}}

% Inserted by MH to get step number etc. to work nicely
% My attempt to define a worked example environment
\newcounter{stepcounter}
\setcounter{stepcounter}{1}

\newcounter{workedexamplecounter}
\setcounter{workedexamplecounter}{1}


%EZ new counters for exercises, 
\newcounter{exercisecounter}
\setcounter{exercisecounter}{1}
% make it reset for new chapters
\makeatletter
\@addtoreset{exercisecounter}{chapter}
\makeatother


% can we deprecate the use of this, and use \westep instead
%\newcommand{\westep}[1]
%{\par \par%\pagebreak[1]
% \textbf{\textit{Step \arabic{stepcounter} : #1}}\par
% \addtocounter{stepcounter}{1}}


\newenvironment{oldwex}[3]
{% This is a picture at the beginning, a line or bar or something else
\begin{quotation}
\noindent
\vspace{.5cm}
\setlength{\parindent}{0pt}
%% Title and counter
\rput(-2.0,0){\pspencil}\noindent
%\begin{mdframed}[linewidth =2,margin=40]
\textbf{\textit{Worked Example \arabic{workedexamplecounter}: }
#1}\par
\textbf{Question:} #2\par
\textbf{Answer}\par #3%
\addtocounter{workedexamplecounter}{1}\setcounter{stepcounter}{1}%
\par%\vspace{.25cm}
}%
{%
%\end{mdframed}
\end{quotation}
\begin{flushright}
\end{flushright}
\setlength{\parindent}{\saveparindentlength}}



\newskip\Vspace@topskip
\providecommand\Vspace[1]{%
        \Vspace@topskip\topskip
        \topskip\z@skip
        \ifvmode
                \dimen@\prevdepth
                \hrule\@height\z@
                \nobreak
                \vskip\glueexpr#1\relax
                \vskip\z@skip
                \prevdepth\dimen@
        \else
                \@bsphack
                \vadjust{%
                        \@restorepar
                        \hrule\@height\z@
                        \nobreak
                        \vskip#1%
                        \vskip\z@skip
                }%
                \@esphack
        \fi
        \topskip\Vspace@topskip
}







% can we deprecate the use of this, and use \westep instead
\newcommand{\westep}[1]
{\par \vspace{3mm}%\pagebreak[1]
 \textit{Step \arabic{stepcounter}} : \textit{\textbf{#1}}\par\vspace{3mm}
 \addtocounter{stepcounter}{1}}

\definecolor{siyavula-gray}{RGB}{239,239,239}
\definecolor{siyavula-dark-gray}{RGB}{194,194,194}

\newenvironment{wex}[3]{
\newlength{\myfloatsep}
\setlength{\myfloatsep}{\floatsep}
\setlength{\floatsep}{2mm}
\nopagebreak 
\begin{quotation}
\noindent
\vspace{.5cm}
\setlength{\parindent}{0pt}
%% Title and counter
%\rput(-2.0,0){\pspencil}\noindent

\begin{mdframed}[linewidth=1.5pt,linecolor=darkgray,innertopmargin=10pt, innerbottommargin=10pt]
{\fontsize{10}\bfseries\slshape Worked Example \arabic{workedexamplecounter}: }#1
\noindent
\end{mdframed}
%\ifdimless{\pagetotal}{4\baselineskip}%4em play with this
%{\Vspace{8\baselineskip}}{}
\begin{mdframed}[skipabove=5mm, skipbelow=5mm, backgroundcolor=siyavula-gray, linecolor=siyavula-gray]
%%% EZ trying a tikzpicture instead of mdframed for the heading. 
%\begin{tikzpicture}[remember picture,overlay]
%    \node[anchor=west, yshift=8pt, xshift=-0.82cm, text centered]
%      {\begin{tikzpicture}[remember picture, overlay]
%        \draw[color=white, draw=black, very thick] (0,0) rectangle
%          (\textwidth,0.75cm);
%       \node[anchor=west, yshift=0.3525cm]
%              {\color{black}{\fontsize{10}]\bfseries\slshape Worked Example \arabic{workedexamplecounter}: } #1 };
%       \end{tikzpicture}
%    };
%\end{tikzpicture}

\vspace*{\baselineskip}
\textbf{\textit{QUESTION}} \vspace*{\baselineskip} 
\hrule
\vspace*{\baselineskip}
\textit{\noindent #2}
\vspace*{2\baselineskip}
\par
\textbf{\textit{SOLUTION}}\vspace*{\baselineskip}
\noindent
\hrule
\vspace*{\baselineskip}
 #3%
\addtocounter{workedexamplecounter}{1}\setcounter{stepcounter}{1}%
\par%\vspace{.25cm}
}%
{%
\end{mdframed}
\end{quotation}
%\nopagebreak[4]
%\begin{pspicture}(0,0)(10,0.01)
% \psline[linecolor=gray,linewidth=3pt]{-}(0,0)(10,0)
% \psline[linecolor=gray]{-}(10,0)(10,2)
%\end{pspicture}
\setlength{\parindent}{\saveparindentlength}
\par

\setlength{\floatsep}{\myfloatsep}

}








\newenvironment{wex2}[3]
{% This is a picture at the beginning, a line or bar or something else
% this command puts a bar in the margin and sets the bar width
% Now we set margins
% \begin{center}
%\pagebreak
%\begin{flushleft}
%\begin{pspicture}(0,0)(10,0.01)
% \psline[linecolor=gray,linewidth=3pt]{-}(0,0)(10,0)
% \psline[linecolor=gray]{-}(0,0)(0,-2)
%\end{pspicture}
%\end{flushleft}
%\end{center}
\begin{quotation}
\noindent
\vspace{./5cm}
\setlength{\parindent}{0pt}
%% Title and counter
\rput(-2.0,0){\pspencil}\noindent
\begin{mdframed}[linewidth=2,linecolor=black,backgroundcolor=lightgray]
\textbf{\textit{Worked Example \arabic{workedexamplecounter}: }
#1}\par
\textbf{Question:} #2\par
\textbf{Answer}\par #3%
\addtocounter{workedexamplecounter}{1}\setcounter{stepcounter}{1}%
\par%\vspace{.25cm}
}%
{%
\end{mdframed}
\end{quotation}
%\nopagebreak[4]
\begin{flushright}
%\begin{pspicture}(0,0)(10,0.01)
% \psline[linecolor=gray,linewidth=3pt]{-}(0,0)(10,0)
% \psline[linecolor=gray]{-}(10,0)(10,2)
%\end{pspicture}
\end{flushright}
\setlength{\parindent}{\saveparindentlength}}



\newenvironment{exercises}[1]
{
\exercisesheading

%\hrule width \textwidth height 1pt\\
\begin{quotation}
\setlength{\parindent}{0pt}
%\vspace{\baselineskip}
%\textbf{\textit{Exercises:}
%#1}
\par
}%
{%
\end{quotation}
\setlength{\parindent}{\saveparindentlength}

%insert the more practise at www.maths10.
( More practise at \url{www.maths10.co.za} with \raisebox{-5 pt}{\includegraphics[width=0.5cm]{col11306.imgs/summary_www.png}} M049 )
}

\newenvironment{eocexercises}[1]
{
\eocexercisesheading
%\hrule width \textwidth height 1pt\\
\begin{quotation}
%\noindent
%\vspace{.5cm}
\setlength{\parindent}{0pt}
%\textbf{\textit{End of chapter exercises:}
%#1}
\par
}%
{%
\end{quotation}
\setlength{\parindent}{\saveparindentlength}}

\newenvironment{f_experiment}[1]
{
\begin{quotation}
\begin{mdframed}[linewidth=5,linecolor=darkgray]
\noindent
\vspace{.5cm}
\setlength{\parindent}{0pt}
\textbf{\textit{Formal Experiment:}
#1}\par
}%
{%
\end{mdframed}
\end{quotation}
\setlength{\parindent}{\saveparindentlength}}


\newenvironment{i_experiment}[1]
{
\begin{quotation}
\begin{mdframed}[linewidth=5,linecolor=lightgray]
\noindent
\vspace{.5cm}
\setlength{\parindent}{0pt}
\textbf{\textit{Informal Experiment:}
#1}\par
}%
{%
\end{mdframed}
\end{quotation}
\setlength{\parindent}{\saveparindentlength}}

\newenvironment{g_experiment}[1]
{
\begin{quotation}
\begin{mdframed}[linewidth=5,linecolor=lightgray]
\noindent
\vspace{.5cm}
\setlength{\parindent}{0pt}
\textbf{\textit{General Experiment:}
#1}\par
}%
{%
\end{mdframed}
\end{quotation}
\setlength{\parindent}{\saveparindentlength}}

\newenvironment{oldactivity}[1]
{
\begin{quotation}
\noindent
\vspace{.5cm}
\setlength{\parindent}{0pt}
\textbf{\textit{Activity:}
#1}\par
}%
{%
\end{quotation}
\setlength{\parindent}{\saveparindentlength}}

\newenvironment{activity}[1]
{
\begin{quotation}
\noindent
\setlength{\parindent}{0pt}
\hrule width \textwidth height 1pt\\
\vspace{1pt}\hrule width \textwidth height 1pt\vspace{7pt}
\textbf{\Large\textit{Activity:}
#1}
\vspace{5pt}
\hrule width \textwidth height 1pt
\vspace{5pt}
}%
{%
\vspace{5pt}
\hrule width \textwidth height 1pt
\end{quotation}
\setlength{\parindent}{\saveparindentlength}}

\newenvironment{groupdiscussion}[1]
{
\begin{quotation}
\noindent
\setlength{\parindent}{0pt}
\hrule width \textwidth height 1pt\\
\vspace{1pt}\hrule width \textwidth height 1pt\vspace{7pt}
\textbf{\Large\textit{Group Discussion:}
#1}
\vspace{5pt}
\hrule width \textwidth height 1pt
\vspace{5pt}
}%
{%
\vspace{5pt}
\hrule width \textwidth height 1pt
\end{quotation}
\setlength{\parindent}{\saveparindentlength}}


\newenvironment{oldgroupdiscussion}[1]
{
\begin{quotation}
\noindent
\vspace{.5cm}
\setlength{\parindent}{0pt}
\textbf{\textit{Group Discussion:}
#1}\par
}%
{%
\end{quotation}
\setlength{\parindent}{\saveparindentlength}}

\newenvironment{casestudy}[1]
{
\begin{quotation}
\noindent
\vspace{.5cm}
\setlength{\parindent}{0pt}
\textbf{\textit{Case Study:}
#1}\par
}%
{%
\end{quotation}
\setlength{\parindent}{\saveparindentlength}}

\newenvironment{project}[1]
{
\begin{quotation}
\noindent
\vspace{.5cm}
\setlength{\parindent}{0pt}
\textbf{\textit{Project:}
#1}\par
}%
{%
\end{quotation}
\setlength{\parindent}{\saveparindentlength}}

\newenvironment{Investigation}[1]
{
\begin{quotation}
\noindent
\vspace{.5cm}
\setlength{\parindent}{0pt}
\textbf{\textit{Investigation:}
#1}\par
}%
{%
\end{quotation}
\setlength{\parindent}{\saveparindentlength}}




        \newenvironment{cnxrule}{\begin{example}}
        {\end{example}}

\newenvironment{exercise}{\textbf{\par Needs to be changed to wex environment\par}\begin{example}}
        {\end{example}}


        \newenvironment{definition}{\begin{example}}
        {\end{example}}

        \newenvironment{listname}{\vspace{\rubberspace}\par\noindent{}}{\vspace*{0pt}
        \nopagebreak}

    \newcommand{\lessthan}{\ensuremath{<}}
    \newcommand{\greatthan}{\ensuremath{>}}

    % --------------------------------------------
    % Hacks for honouring row/entry/@align
    % (\hspace not effective when in paragraph mode)
    % Naming convention for these macros is:
    % 'docbooktolatex' 'align' {alignment-type} {position-within-entry}
    % where r = right, l = left, c = centre
    \newcommand{\docbooktolatexalign}[2]{\protect\ifvmode#1\else\ifx\LT@@tabarray\@undefined#2\else#1\fi\fi}
    \newcommand{\docbooktolatexalignll}{\docbooktolatexalign{\raggedright}{}}
    \newcommand{\docbooktolatexalignlr}{\docbooktolatexalign{}{\hspace*\fill}}
    \newcommand{\docbooktolatexaligncl}{\docbooktolatexalign{\centering}{\hfill}}
    \newcommand{\docbooktolatexaligncr}{\docbooktolatexalign{}{\hspace*\fill}}
    \newcommand{\docbooktolatexalignrl}{\protect\ifvmode\raggedleft\else\hfill\fi}
    \newcommand{\docbooktolatexalignrr}{}

    % ------------------------------------------
    % Break long chapter titles in running heads
    \makeatletter
      \def\@evenhead{\thepage\hfil\parbox{5in}{\raggedleft\slshape\leftmark}}%
    \makeatother

    % ------------------------------------------
    % Lengths etc. for tables
    \newlength\mytablewidth  % full width of table
    \newlength\mytablespace  % non-content hspace in table (tabcolseps and rule widths)
    \newlength\mytableroom   % content hspace in table
    \newlength\mycolwidth
    \newlength\myfixedwidth  % sum of widths of columns that have fixed widths
    \newlength\mystarwidth   % length of one star factor
    \newlength\myspanwidth
    \newsavebox\mytablebox
    \newlength\mytableboxwidth
    \newlength\mytableboxheight
    \newlength\mytableboxdepth
    \newcolumntype{L}[1]{>{\raggedright\hspace{0pt}}p{#1}}
    \newcolumntype{C}[1]{>{\centering\hspace{0pt}}p{#1}}
    \newcolumntype{R}[1]{>{\raggedleft\hspace{0pt}}p{#1}}
    \newcolumntype{J}[1]{>{\hspace{0pt}}p{#1}}

    % ------------------------------------------
    % savebox and length for conditional typesetting of block math
    \newsavebox\mymathbox
    \newlength\mymathboxwidth

    % ------------------------------------------
    % Command to help dmath* be pseudo-text-math
    \newcommand*\nodisplayskips{%
      \setlength\abovedisplayskip{0pt}%
      \setlength\abovedisplayshortskip{0pt}%
      \setlength\belowdisplayskip{0pt}%
      \setlength\belowdisplayshortskip{0pt}%
    }

    \newlength\figurerulewidth
    \setlength\figurerulewidth{\textwidth}
    \addtolength\figurerulewidth{-50pt}
    
    \newlength\rubberspace
    \setlength\rubberspace{3pt plus 3pt}
    
    % Give tables more vertical padding, so that math looks better
    \renewcommand{\arraystretch}{1.4}
    
    
  
    \usepackage{fancyhdr}
    \pagestyle{fancy}
    \fancyfoot[C]{}
    \fancyfoot[LE,RO]{\thepage}



\usepackage[Bjornstrup]{fncychap}
\ChNumVar{\fontsize{76}{80}\usefont{OT1}{pzc}{m}{n}\selectfont}
\ChTitleVar{\raggedleft\HUGE\sffamily\bfseries}



% border icons
\newcommand{\psintfact}%
{%
\rput(0,0){\begin{tabular}{c}\textbf{Interesting}\\ \textbf{Fact}%
\end{tabular}}
}

\newcommand{\psfact}%
{%
\begin{pspicture}(-.5,-.5)(.5,.5)
 \psintfact
 \PstLens[LensSize=0.5,LensRotation=-60,LensMagnification=1.5](-.04,-.1){\psintfact}
\end{pspicture}
}





%Physics Definitions
\def\rope{\psline(0,0.2)(5,.2)\psline(0,0)(5,0)\multirput(0,0)(0.2,0){25}{\psline(0,0)(0.2,0.2)}}

\def\cart{\psframe(0,0.4)(1,1)\pscircle(0.2,0.2){0.2}\pscircle(0.8,0.2){0.2}}

% Arrow for objects and images
\newpsobject{oi}{psline}{arrowsize=6pt, arrowlength=1.5, arrowinset=0, linewidth=2pt}
\psset{lensHeight=3,lensColor=lightgray}
\newpsobject{PrincipalAxis}{psline}{linewidth=0.5pt,linecolor=gray}
\newcommand{\ms}{m$\cdot$s$^{-1}$}                %m/s in text







% fancy headings

\usepackage{tikz}
\usepackage[explicit]{titlesec}

\newcommand*\chapterlabel{}

\titleformat{\chapter}
  {\gdef\chapterlabel{}
   \normalfont\sffamily\fontsize{32}\bfseries\slshape}
  {\gdef\chapterlabel{\thechapter\ }}{0pt}
  {\begin{tikzpicture}[remember picture,overlay]
    \node[yshift=-3cm] at (current page.north west)
      {\begin{tikzpicture}[remember picture, overlay]
        \shade[left color=white, right color=gray] (0,0) rectangle
          (1.2\paperwidth,3cm);
        \node[anchor=west,text centered, yshift=1.5cm, xshift=0.2\paperwidth]
              {\color{black} #1 \hfill \color{white}\chapterlabel};
       \end{tikzpicture}
      };
   \end{tikzpicture}
  }
\titlespacing*{\chapter}{1.5cm}{1cm}{1cm}



\newcommand*\sectionlabel{}
\titleformat{\section}
  {\gdef\sectionlabel{}
   \normalfont\sffamily\fontsize{18}\bfseries\slshape}
  {\gdef\sectionlabel{\thesection }}{0pt}
  {\begin{tikzpicture}[remember picture,overlay]
    \node[anchor=west,yshift=-1.5cm, xshift=-0.4\paperwidth]
      {\begin{tikzpicture}[remember picture, overlay]
        \draw[fill=siyavula-gray, color=siyavula-gray] (0,0) rectangle
          (1.3\paperwidth,1cm);
        \node[anchor=west,text centered ,yshift=0.5cm, xshift=0.3\paperwidth]
              {\color{black} #1  };
       \end{tikzpicture}
      };
   \end{tikzpicture}
  }
\titlespacing*{\section}{1.5cm}{1.5cm}{2cm}



\newcommand*\subsectionlabel{}
\titleformat{\subsection}
  {\gdef\subsectionlabel{}
   \normalfont\sffamily\fontsize{18}\bfseries\slshape}
  {\gdef\subsectionlabel{\thesubsection\ }}{0pt}
  {\begin{tikzpicture}[remember picture,overlay]
   \node[anchor=west, yshift=-1.5cm, xshift=-0.085\paperwidth]
      {\begin{tikzpicture}[remember picture, overlay]
        \draw[fill=siyavula-gray, color=siyavula-gray] (0,0) rectangle
          (0.915\paperwidth,1.0cm);
        \node[anchor=west,text centered ,yshift=0.5cm]
              {\color{black} #1  };
       \end{tikzpicture}
      };
   \end{tikzpicture}
  }
\titlespacing*{\subsection}{1.5cm}{1cm}{2cm}



\newcommand{\summary}{
\vspace{6\baselineskip} \par
{\normalfont \sffamily\fontsize{22}\bfseries
\begin{tikzpicture}[remember picture,overlay]
    \node[anchor=west, xshift=-0.4\paperwidth]
      {\begin{tikzpicture}[remember picture, overlay]
        \draw[fill=siyavula-gray, color=siyavula-gray] (0,0) rectangle
          (1.2\paperwidth,2cm);
        \node[anchor=west,text centered ,yshift=1.cm, xshift=0.3\paperwidth]
              {   \color{black} Chapter \chapterlabel | Summary  };
       \end{tikzpicture}
      };
   \end{tikzpicture}

}
\par
\vspace{\baselineskip}
}


\newcommand{\exercisesheading}{
\addtocounter{exercisecounter}{1}
 \par  \vspace*{3\baselineskip}
{  \fontsize{10}\bfseries\itshape
\begin{tikzpicture}[remember picture,overlay]
    \node[anchor=west ,text centered]
      {\begin{tikzpicture}[remember picture, overlay]
        \draw[color=white, draw=black, very thick] (0,0) rectangle
          (0.33\textwidth, 0.75cm);
       \node[anchor=west, yshift=0.3525cm, text centered]
              { \bfseries ~ \color{black} Exercise \chapterlabel~-~\arabic{exercisecounter}   };
       \end{tikzpicture}
    };
    
% now the gray bar to the end of the paragraph
    \node[anchor=west, xshift=0.335\textwidth]{
      \begin{tikzpicture}[remember picture, overlay]
        \draw[fill=siyavula-gray, color=siyavula-gray] (0,0) rectangle
          (0.655\textwidth,0.75cm);
       \end{tikzpicture}


      };
   \end{tikzpicture}

}
\par
\vspace{\baselineskip}

}
\newcommand{\eocexercisesheading}{
\par \vspace*{3\baselineskip} 
{\normalfont \sffamily\fontsize{16}\bfseries
\begin{tikzpicture}[remember picture,overlay]
    \node[anchor=west, yshift=6\baselineskip, xshift=0., text centered]
      {\begin{tikzpicture}[remember picture, overlay]
        \draw[color=white, draw=black, very thick] (0,0) rectangle
          (0.3\textwidth,1.5cm);
       \node[anchor=west, yshift=0.75cm, text centered]
              { ~~  \color{black} Chapter \chapterlabel };
       \end{tikzpicture}
    };
    
% now the gray bar to the end of the paragraph
    \node[anchor=west, xshift=0.31\textwidth]{
      \begin{tikzpicture}[remember picture, overlay]
        \draw[fill=siyavula-gray, color=siyavula-gray] (0,0) rectangle
          (0.69\textwidth,1.5cm);
            \node[anchor=west, yshift=0.75cm, text centered]
              {   \color{black} End of Chapter Exercises };

       \end{tikzpicture}


      };
   \end{tikzpicture}

}
\par
\vspace{\baselineskip}
}

